% generated by roam.py handbook export -- не редактировать, править заметки
% 20261004191857-дихотомия_по_производнои.org (body 4567407f7069)

\begin{method}[{Дихотомия по производной}]\label[method]{note:124BD9BC-084E-43FB-B617-AF8690B33EF8}
  Пусть $f$  -- дифференцируема на $[a;b]$ и \hyperref[note:CD52D818-AA25-471C-97C2-5DBE9787CE53]{унимодальна}.

  Считаем также, что $f^{\prime}(a) < 0$ (минимум лежит правее), $f^{\prime}(b) > 0$ (минимум лежит левее)

  Идея: искать корень производной (по сути бин. поиском).
  Бин. поиск будет работать так как производная $<0$ до минимума и $> 0$ после.
  Точнее: $f^{\prime}(x) > 0 \Rightarrow x^{*} < x$, а $f^{\prime}(x) < 0 \Rightarrow x^{*} > x$.

  Пытаемся найти бин. поиском ту точку, где $\mathrm{sign} f^{\prime}(x) = 0$:
  Берем $x = \frac{a+b}{2}$. Если $f^{\prime}(x) > 0 \Rightarrow b = x$, иначе $a = x$

  Каждый раз отрезок сокращаем вдвое. Тогда \hyperref[note:C42D2C19-B35F-4C0E-AFEB-D880CBB72176]{невязка}
  $r_k= \|x_k - x^{*}\| \le \frac{b-a}{2^k}$. Скорость сходимости~--- \hyperref[note:DF2405BF-8B31-47EB-8DEF-B1A88034717C]{линейная}
\end{method}
